% Generated draft: requires mathematical and source review.
\documentclass[11pt]{article}
\usepackage[margin=25mm]{geometry}
\usepackage{fontspec}
\setmainfont{Latin Modern Roman}
\usepackage{amsmath,amssymb,mathtools,unicode-math}
\setmathfont{Latin Modern Math}
\usepackage{enumitem,needspace,xurl,hyperref}
\setlength{\parindent}{0pt}
\setlength{\parskip}{5pt}
\setlength{\emergencystretch}{3em}
\newcommand{\N}{\mathbb{N}}
\newcommand{\Z}{\mathbb{Z}}
\newcommand{\Q}{\mathbb{Q}}
\newcommand{\R}{\mathbb{R}}
\newcommand{\C}{\mathbb{C}}
\newcommand{\F}{\mathbb{F}}
\newcommand{\A}{\mathbb{A}}
\newcommand{\PP}{\mathbb{P}}
\newcommand{\E}{\mathbb{E}}
\newcommand{\eps}{\varepsilon}
\newcommand{\dd}{\,\mathrm{d}}
\begin{document}
\section*{UnsolvedMath ID: 700738}
\textbf{Source problem code:} catalogue entry
\subsection{Definitions and mathematical statement}
For a group \(G\), write \(1\) for its identity, set \([x,y]=x^{-1}y^{-1}xy\), and use left-normed iterated commutators, so that \([x,y,y]=[[x,y],y]\). For subgroups \(H,K\leq G\), let \([H,K]\) be the subgroup generated by their commutators. Define \(\gamma_1(G)=G\) and \(\gamma_{i+1}(G)=[\gamma_i(G),G]\). The group \(G\) has nilpotency class exactly \(c\) if \(\gamma_{c+1}(G)=\{1\}\) and \(\gamma_c(G)\neq\{1\}\). A finite \(p\)-group is a group of order \(p^n\), where \(p\) is prime and \(n\) is a nonnegative integer.

Let \(\operatorname{End}(G)\) denote the endomorphisms of \(G\), and put \(a^\varphi=\varphi(a)\). Let \(\operatorname{Aut}(G)\) denote its automorphism group, let \(Z(G)\) denote its center, and define
\[
\operatorname{Aut}_c(G)=\{\alpha\in\operatorname{Aut}(G):x^{-1}\alpha(x)\in Z(G)\text{ for every }x\in G\}.
\]
Thus these are the central automorphisms. Let \(\operatorname{Inn}(G)\) be the subgroup consisting of the automorphisms \(x\mapsto g^{-1}xg\), with \(g\in G\). The equality \(\operatorname{Aut}(G)=\operatorname{Aut}_c(G)\operatorname{Inn}(G)\) means that every automorphism is a composition of a central automorphism and an inner automorphism. A group is \(2\)-Engel if \([x,y,y]=1\) for every \(x,y\in G\).

\begin{enumerate}
\item[(a)] Does there exist a finite \(3\)-group \(G\) of nilpotency class exactly \(3\) such that
\[
\forall a\in G\ \forall\varphi\in\operatorname{End}(G),\qquad [a,\varphi(a)]=1?
\]
\item[(b)] For a prime \(p\), consider the existence assertion
\[
\exists\,G\text{ a finite }p\text{-group}:\quad
\operatorname{class}(G)>2
\quad\text{and}\quad
\operatorname{Aut}(G)=\operatorname{Aut}_c(G)\operatorname{Inn}(G).
\]
This is the existence question in the supplied archived part (b).
\item[(c)] For a prime \(p\), consider the existence assertion
\[
\exists\,G\text{ a finite }p\text{-group}:\quad
\operatorname{class}(G)>2,\qquad
\forall x,y\in G,\ [x,y,y]=1,
\qquad
\operatorname{Aut}(G)=\operatorname{Aut}_c(G)\operatorname{Inn}(G).
\]
This is the existence question in the supplied archived part (c).
\end{enumerate}

The source does not prescribe a value of \(p\) or explicitly quantify over primes in (b) and (c). Their assertions have therefore been displayed with \(p\) as a prime parameter; no claim that they hold for every prime is intended. The supplied archive reports affirmative answers to (b) and (c), attributed respectively to I. Malinowska and to A. Abdollahi, A. Faghihi, S. A. Linton, and E. A. O'Brien. The supplied registry labels the overall problem partially solved and retains (a) in the main section. These are supplied source assessments, not independently verified mathematical conclusions.
\subsection{Short English statement}
Part (a) asks for a finite group of order a power of \(3\), of class exactly \(3\), in which every element commutes with its image under every endomorphism. Archived parts (b) and (c) ask for finite \(p\)-groups of class greater than \(2\) whose automorphisms are products of central and inner automorphisms; part (c) also requires the \(2\)-Engel identity. The supplied archive reports affirmative answers to (b) and (c). The intended quantification over primes in those parts is not explicit.
\subsection{Sources}
\begin{itemize}
\item Supplied local UnsolvedMath registry record: permanent ID \(700738\), source problem code \texttt{catalogue entry}, title \emph{Kourovka Notebook Problem 11.46}, upstream revision \texttt{d371ee917ac9647fbb06cd6fc529deec6f3dda82}. The supplied status is \texttt{partially\_solved}. Its status review is dated \texttt{2026-10-16}; this assessment is supplied dated curation, not independently verified.
\item \emph{The Kourovka Notebook (Unsolved Problems in Group Theory)}, 11th issue (1990), Problem 11.46, proposed by A. Caranti; supplied record reports an update in the 21st issue (2026), \texttt{arXiv:1401.0300v48}. The supplied collection attribution names E. I. Khukhro and V. D. Mazurov as editors.
\item Reference supplied with part (a): A. Caranti, \emph{J. Algebra} \textbf{97}, no. 1 (1985), 1--13. No article title or link is supplied.
\item Archive reference supplied for the affirmative answer to part (b): I. Malinowska, \emph{Rend. Sem. Mat. Padova} \textbf{91} (1994), 265--271. No article title or link is supplied.
\item Archive reference supplied for the affirmative answer to part (c): A. Abdollahi, A. Faghihi, S. A. Linton, and E. A. O'Brien, \emph{Arch. Math. (Basel)} \textbf{95}, no. 1 (2010), 1--7. No article title or link is supplied. The complete bibliographic details occur in the supplied background; the status-review excerpt is truncated.
\end{itemize}
\end{document}
